\chapter{Estrutura do projeto}

Voce ja tem uma estrutura inicial. A ideia e manter essa base e agregar os modulos necessarios aos poucos.

\section{Estrutura inicial fornecida}

\begin{verbatim}
apple_health_export
├─ .editorconfig
├─ README.md
├─ notebooks
│  ├─ README.md
│  └─ playground.ipynb
├─ poetry.lock
├─ pyproject.toml
└─ src
   ├─ __init__.py
   ├─ app.py
   ├─ config
   │  ├─ __init__.py
   │  ├─ logger.py
   │  └─ settings.py
   └─ parser
      ├─ __init__.py
      ├─ base_parser.py
      └─ xml_parser.py
\end{verbatim}

Essa estrutura ja tem tres decisoes boas:

\begin{itemize}
  \item \code{src/app.py}: ponto de entrada;
  \item \code{src/config}: configuracao e logging;
  \item \code{src/parser}: responsabilidade de parsing separada.
\end{itemize}

\section{Estrutura agregada proposta}

Agora vamos expandir sem destruir sua organizacao:

\begin{verbatim}
apple_health_export
├─ .editorconfig
├─ README.md
├─ docker-compose.yml
├─ .env.example
├─ notebooks
│  ├─ README.md
│  └─ playground.ipynb
├─ poetry.lock
├─ pyproject.toml
├─ tests
│  ├─ __init__.py
│  ├─ test_parser_utils.py
│  ├─ test_xml_parser.py
│  └─ fixtures
│     └─ small_export.xml
└─ src
   ├─ __init__.py
   ├─ app.py
   ├─ config
   │  ├─ __init__.py
   │  ├─ logger.py
   │  └─ settings.py
   ├─ db
   │  ├─ __init__.py
   │  ├─ connection.py
   │  ├─ schema.py
   │  └─ repository.py
   ├─ models
   │  ├─ __init__.py
   │  └─ health_record.py
   ├─ parser
   │  ├─ __init__.py
   │  ├─ base_parser.py
   │  ├─ xml_parser.py
   │  ├─ gpx_parser.py
   │  └─ profiler.py
   ├─ pipelines
   │  ├─ __init__.py
   │  ├─ sync_pipeline.py
   │  ├─ threaded_pipeline.py
   │  └─ multiprocessing_pipeline.py
   └─ insights
      ├─ __init__.py
      ├─ daily_summary.py
      └─ queries.py
\end{verbatim}

\section{Por que cada pasta existe?}

\begin{longtable}{p{4cm}p{9cm}}
\toprule
\textbf{Pasta/arquivo} & \textbf{Responsabilidade} \\
\midrule
\code{config/settings.py} & centralizar variaveis como host do banco, porta, usuario, senha, batch size \\
\code{config/logger.py} & configurar logs padronizados \\
\code{models/} & dataclasses e objetos de dominio \\
\code{parser/} & transformar XML/GPX em objetos Python \\
\code{db/} & conexao, schema e escrita no Postgres \\
\code{pipelines/} & orquestrar parser, fila, batches e writer \\
\code{insights/} & queries e agregacoes para gerar metricas \\
\code{tests/} & testes pequenos sem usar o XML gigante \\
\bottomrule
\end{longtable}

\section{Regra de separacao}

Cada modulo deve ter uma responsabilidade clara:

\begin{itemize}
  \item Parser nao deve saber detalhes do Docker.
  \item Repository nao deve parsear XML.
  \item Logger nao deve conter regra de negocio.
  \item Pipeline orquestra componentes, mas nao deve concentrar toda logica.
\end{itemize}

Essa separacao facilita testes e torna a evolucao mais segura.
